% titel: Blatt 6, A2: Rechnen mit komplexen Zahlen
% skript: Kap. 1 (Komplexe Zahlen, Betrag und Abstand)

% Teil I: Werkzeuge, kein Lösungsweg.

Ungleichungen mit dem Betrag in $\mathbb C$ lesen, konkrete Lösungen
finden und die Multiplikativität des Betrags zeigen.

\begin{werkzeugkasten}[Betrag als Abstand -- und was daraus \emph{gezeichnet} wird]
Für $z=a+ib$ ist $|z|=\sqrt{a^2+b^2}$ der euklidische Abstand zum
Ursprung; $|z-z_0|$ ist der Abstand zwischen $z$ und $z_0$. Deshalb werden
Betragsgleichungen und -ungleichungen zu \emph{geometrischen} Aussagen:
eine Forderung an den Abstand zu einem festen Punkt. Beispiel:
\[
  |z-1+i|=|z-(1-i)|>1
\]
heißt: alles, was weiter als $1$ vom Punkt $1-i$ entfernt ist -- das
Außengebiet des Einheitskreises um $1-i$, ohne Rand. Das ist der Trick:
$\pm i$ als Punkte in der Ebene lesen, nicht als \emph{Zahlen für sich}.

\textit{Skript:} Kap. 1 (Betrag).
\end{werkzeugkasten}

\begin{werkzeugkasten}[Schnitt zweier Kreise -- systematisch]
Um alle $z$ mit $|z-z_1|=r_1$ und $|z-z_2|=r_2$ zu finden, nutzt man,
dass Quadrieren die Beträge in quadratische Koordinaten-Ausdrücke
verwandelt:
\begin{enumerate}
  \item Beide Gleichungen quadrieren:
        $|x+iy-z_0|^2=(\cdots)^2+(\cdots)^2=r^2$.
  \item \textbf{Subtrahieren} einer von der anderen. Die
        $x^2$-Terme fallen weg, es bleibt eine \emph{lineare} Beziehung
        zwischen $x$ und $y$ (die Radikal-Linie/Sekante).
  \item Linear-Nichtlineares System lösen: eine Variable eliminieren, die
        quadratische Restgleichung lösen. Zwei Kreise schneiden sich in
        \emph{zwei} Punkten (oder einem, oder keinem).
\end{enumerate}
Für die allgemeine Aussage $|zw|=|z||w|$ rechnet man algebraisch mit
$|u|^2=u\bar u$ und der Rechenregel $\overline{zw}=\bar z\,\bar w$ des
Konjugierens.

\textit{Skript:} Kap. 1 (Betrag, Konjugation).
\end{werkzeugkasten}

\begin{achtung}[Typische Fallen]
\begin{itemize}
  \item \textbf{Betrag ist nicht additiv:} $|z-1+i|$ ist \emph{nicht}
        $|z-1|+|i|$; der Betrag wirkt auf die gesamte Differenz.
  \item \textbf{Quadratieren verliert nichts:} Alle Terme sind quadratisch,
        es geht kein Vorzeichen \emph{verloren} (anders als beim Quadrieren
        einer Gleichung mit linearen Termen).
  \item \textbf{Alle Schnittpunkte angeben:} Zwei Kreise liefern in der
        Regel \emph{zwei} Lösungen, nicht eine.
  \item \textbf{$|u|^2=u\bar u$ ist der saubere Weg} für die
        Multiplikativität; direkt \emph{ausmultiplizieren} von $|zw|$ macht die
        Sache unübersichtlich.
\end{itemize}
\end{achtung}
